\documentclass[11pt]{article}
\usepackage[margin=1in]{geometry}
\usepackage[T1]{fontenc}
\usepackage{lmodern,amsmath,amssymb,amsthm,booktabs,microtype,graphicx}
\usepackage[colorlinks=true,linkcolor=blue,urlcolor=blue,citecolor=blue]{hyperref}

\newtheorem{theorem}{Theorem}[section]
\newtheorem{lemma}[theorem]{Lemma}
\newtheorem{proposition}[theorem]{Proposition}
\newtheorem{corollary}[theorem]{Corollary}
\theoremstyle{remark}
\newtheorem{remark}[theorem]{Remark}

\newcommand{\E}{\mathbb E}
\newcommand{\Prb}{\mathbb P}
\newcommand{\one}{\mathbf 1}
\newcommand{\diag}{\operatorname{diag}}
\newcommand{\eps}{\varepsilon}
\newcommand{\Cov}{\operatorname{Cov}}
\newcommand{\Id}{\mathrm{Id}}

\title{Inherited-family dependence can reverse an\\ equal-exposure treatment decision}
\author{}
\date{21 September 2026}

\begin{document}
\maketitle

\begin{abstract}
We exhibit two seven-type branching models that agree on every type-resolved
population mean, under every common deterministic schedule, and nevertheless
prefer opposite orders for two equal-exposure interventions. The models differ
only in the joint law of sister birth states: one partitions the mother's
epigenetic marks complementarily between daughters, the other draws both
daughters independently from the very same one-daughter marginal. All division,
death, switching and mutation rates, the one-mutant-daughter convention and
both exposures are identical. The objective is eventual total extinction under
a declared common continuation, so the comparison is a probability, not a mean.
Exact outward arithmetic places the four one-founder extinction probabilities in
rational intervals of width $3\times10^{-15}$ and proves opposite order margins
exceeding $1.6\times10^{-6}$; an analytic curvature bound together with four
additional enclosed flows certifies the reversal throughout an interval of
killing rates of width $1.6\times10^{-5}$. We identify the mechanism: the two
laws share their first-order response to rare acquisition, and their second-order
responses differ by the transported sister covariance. We prove this correction
has a definite sign, so the independent approximation is systematically
optimistic about eradication, and we give an exact finite-amplitude response
identity valid at the witness parameters. The operational cost of the
approximation is the loss $1.7013\times10^{-6}$ in extinction probability
incurred by following its recommendation when the complementary law holds; a
decision-preservation criterion says when a schedule ranking survives the
approximation, and a scalar survival floor bounds what any predictable feedback
course can achieve. Two continuation perturbations are quantified separately:
residual acquisition up to $3\times10^{-5}$ and finite clearing of duration
$152$, the latter at a common exposure cost $15.2$ times the course itself.
Arguments are conventional probability and ODE proofs, numerical claims are
exact interval arithmetic under a written soundness proof, and a thin Lean root
checks selected algebra but not the stochastic theorem. This is a synthetic
counterexample with a small near-tie effect, not a calibrated treatment
recommendation.
\end{abstract}

\section{Introduction}\label{sec:intro}

\subsection{A model-reduction question}
Reversible tolerance can sustain a cell population long enough for durable
resistance to arise, and branching models already connect phenotypic switching,
permanent resistance and drug sequencing \cite{gunnarsson2020,iwasa2006}.
Dose-response and induction laws shape optimal dosing \cite{gunnarsson2025}, and
the timing of a second strike has been analysed directly in terms of
evolutionary rescue \cite{patil2026,orr2014}. Neither treatment sequencing nor
the general observation that offspring dependence can affect extinction is new
here.

Our question is narrower and is about model reduction. When a stochastic
population model is simplified, a common and reasonable defence is that the
simplification preserves the observable mean dynamics. Suppose it does so
exactly: fix the molecular dynamics, all demographic rates, the acquisition
convention, every one-daughter marginal, the two actions and their exposures,
and change only the \emph{dependence between the two sister states produced at a
division}. Every type-resolved mean trajectory is then unchanged, under every
common schedule. Can the better treatment order for an exact extinction
objective nevertheless change?

\subsection{What we prove}
It can. We construct a literal seven-type source (Section~\ref{sec:source}) with
six molecular types carrying activating and repressive marks and one durable
resistant type $M$. Under the \emph{complementary} law $J$ each old mark of the
mother is allocated to one daughter or the other by a fair coin, so the pair
conserves the mother's marks. Under the \emph{independent} law $I$ both daughters
are drawn independently from exactly the marginal that $J$ induces. The two laws
have identical marginals and, as a consequence of the one-mutant-daughter
convention, identical complete mean operators (Proposition~\ref{prop:means}).

Two actions $A$ and $B$ of ten time units each are compared in both chronological
orders, with equal deadline and equal integrated rate increments. The objective
is eventual total extinction under a declared post-course continuation which we
prove clears every finite molecular population almost surely and creates no new
$M$ (Proposition~\ref{prop:endpoint}). Theorem~\ref{thm:decision} states the
counterexample: from any number of identical $AA$ founders, $J$ strictly prefers
$BA$ and $I$ strictly prefers $AB$, with one-founder margins above
$1.6\times10^{-6}$, and both preferences persist throughout a certified interval
of the killing rate in $B$.

Section~\ref{sec:mechanism} explains why. Expanding in the acquisition scale
$\eps$, the two laws share the first-order coefficient exactly; the second-order
coefficients differ by the transported covariance between sisters
(equation~\eqref{eq:covtransport}). We prove in Theorem~\ref{thm:sign} that this
covariance is nonpositive coordinatewise, so the independent approximation
over-estimates extinction probability at leading order --- it is systematically
optimistic about eradication. That sign alone does not decide a schedule
comparison, because a reversal depends on the \emph{difference} of two such
errors; Proposition~\ref{prop:exact} supplies an exact finite-amplitude identity
for that difference, and the schedule-gap identity~\eqref{eq:gapidentity}
connects it to the certified numbers.

Section~\ref{sec:decision} gives the operational content: the direct cost of
acting on the approximation, a criterion under which a schedule ranking is
protected from it, the founder count at which the advantage is largest, and a
minimax reading of the information limitation. Section~\ref{sec:floor} proves a
survival floor valid for \emph{every} admitted predictable feedback course, so
the exhibited schedules are compared against a genuine feasibility limit.
Section~\ref{sec:continuation} quantifies how much the declared continuation can
be perturbed before the conclusion is lost, and states the exposure cost of the
finite version honestly.

\subsection{What is not claimed}
We do not optimise over schedules, identify the maximal reversal set, or refute
any model-specific timing theorem proved under different assumptions. The effect
is small and sits near a tie by construction: that is the point, since a decision
procedure is exactly where a small modelling error becomes consequential. All
rates are synthetic mathematical assumptions, not calibrated drug
concentrations, doses or clinical cure probabilities, and nothing here is a
treatment recommendation. The acquisition probabilities are stipulated
coarse-grained per-division quantities and are not per-base mutation rates.

\subsection{Evidence conventions}\label{sec:evidence}
Claims carry one of four labels, used consistently and summarised in
Appendix~\ref{sec:evidencemap}. \textbf{C} is a conventional written proof.
\textbf{E} is a computation in exact rational or outward dyadic interval
arithmetic, rigorous under a conventional soundness proof given in
Appendix~\ref{sec:certificate}; the arithmetic is exact, but the integrator
itself is not machine-verified. \textbf{K} is a strictly compiled Lean module;
these cover selected algebra only and no Lean theorem asserts the
probability-to-ODE bridge or the complete result. \textbf{N} is a floating-point
diagnostic, used for illustration and never as an input to a certified sign.

\section{The literal source and the extinction endpoint}\label{sec:source}

\subsection{Molecular types, marks and rates}
A molecular cell carries two sites, each unmarked, activating or repressive. Its
type is the pair $(a,r)$ of activating and repressive mark counts with
$a+r\le2$; there are six such types, and we write $u=2-a-r$ for the number of
unmarked sites.

\begin{center}
\begin{tabular}{lrrrrrr}\toprule
Type&$UU$&$UR$&$RR$&$AU$&$AR$&$AA$\\\midrule
Activating marks $a$&0&0&0&1&1&2\\
Repressive marks $r$&0&1&2&0&1&0\\\bottomrule
\end{tabular}
\end{center}

Throughout, $r$ denotes a repressive mark count. The extinction probability of
the resistant family is written $\rho$, never $r$. The nonzero molecular jump
rates are
\begin{align}
 (a,r)\to(a+1,r)&:u(.01+a/2),&
 (a,r)\to(a,r+1)&:u(.01+r/2),\label{eq:source}\\
 (a,r)\to(a-1,r)&:a(e+r/4),&
 (a,r)\to(a,r-1)&:r(e+a/4).\label{eq:source2}
\end{align}
Writing and erasing are each autocatalytic in their own mark and antagonised by
the opposite mark. Vanishing prefactors suppress invalid boundary transitions.
Let $Q_e$ be the resulting conservative generator, with diagonal minus the
outgoing sum. Division has constant rate $b=.1$; molecular death is
$d_i=d_i^0+c$, where $d_i^0=.01$ for the protected types $a>r$, namely $AU$ and
$AA$, and $d_i^0=.3$ otherwise. The erasure level $e$ and the additional death
$c$ are the two action variables.

The seventh type $M$ is durable resistance. It divides into two $M$ cells at rate
$\beta=.1$, dies at rate $\delta=.02$, never reverts, and is unaffected by either
action. Its stipulated rates do not depend on the mark counts of the cell it came
from. Its extinction probability from one cell is the minimal first-event fixed
point
\begin{equation}
 \rho=\delta/\beta=1/5.\label{eq:rho}
\end{equation}
Note that $RR$ means ``two repressive marks'' and is a molecular type; it is not
the resistant type $M$.

\subsection{Two sister laws with the same marginals}
Under the complementary law $J$, each of the mother's $a$ activating and $r$
repressive marks is assigned independently to one of the two daughters by a fair
coin:
\begin{equation}
 (Y,Z)=\bigl((I,J_r),\,(a-I,r-J_r)\bigr),\qquad
 I\sim\mathrm{Bin}(a,1/2),\quad J_r\sim\mathrm{Bin}(r,1/2),\label{eq:Jlaw}
\end{equation}
independently of each other. Unmarked sites refill each daughter to two sites.
The pair conserves the mother's old marks exactly. For molecular test functions
$f,g$ write
\[
 D_{D,i}(f,g)=\E_{D,i}\bigl[f(Y)g(Z)\bigr],\qquad
 (\Lambda f)_i=\E_i\bigl[f(Y)\bigr].
\]
The comparison law $I$ draws both daughters independently from that same
one-daughter marginal, so
\begin{equation}
 D_{I,i}(f,g)=(\Lambda f)_i(\Lambda g)_i.\label{eq:Ilaw}
\end{equation}
Both laws are exchangeable and both induce the same $\Lambda$, which is the only
sense in which they agree at the level of a single daughter. Law $I$ is an
approximation: it generally does not conserve the mother's old marks across the
pair. It is exactly the reduction obtained by recording the empirical
one-daughter distribution and then assuming sisters are independent.

\subsection{Division-associated acquisition of durable resistance}
At each molecular division, with probability $\mu_i=\eps\kappa_i$, exactly one of
the two sampled daughters, chosen uniformly, is replaced by an $M$ cell;
otherwise both daughters retain their molecular types. We set $\kappa_i=.4$ for
the protected types $AU,AA$ and $\kappa_i=.1$ otherwise, and admit
$0\le\eps\le5/2$, which is exactly the range making every $\mu_i$ a probability.
The mother is removed and exactly two daughters remain. Exchangeability gives
the offspring probability generating function
\begin{equation}
 G_{D,i}(x,y)=(1-\mu_i)D_{D,i}(x,x)+\mu_i\,y\,(\Lambda x)_i.\label{eq:offspring}
\end{equation}
This convention is part of the model statement. Mutating each daughter
independently, or adding a third cell, changes \eqref{eq:offspring} and is a
different source; we do not study those variants. We call the event
\emph{division-associated acquisition of durable resistance}, and emphasise that
it is a stipulated coarse-grained probability per division, not a per-base DNA
mutation probability.

\begin{figure}[ht]
\centering\includegraphics[width=\linewidth]{figures/source.pdf}
\caption{The comparison changes only sister dependence. Both laws use the same
daughter marginals, the same demographic rates and the same one-mutant-daughter
convention. The declared continuation is part of the extinction endpoint.}
\label{fig:source}
\end{figure}

\subsection{Backward field and the shared mean operator}
The backward field on the invariant cube $[0,1]^7$ is
\begin{align}
 \Phi_{D,i}(x,y)&=(Q_ex)_i+d_i(1-x_i)+b\{G_{D,i}(x,y)-x_i\},\label{eq:field}\\
 \Phi_M(y)&=\delta(1-y)+\beta(y^2-y).\label{eq:fieldM}
\end{align}
By \eqref{eq:rho}, $\Phi_M(\rho)=0$: the mutant coordinate is stationary at
$\rho$, for every action, so from the terminal value $\rho$ it stays there and
the molecular block decouples into a six-dimensional autonomous system with
$y\equiv\rho$. We use this repeatedly.

\begin{proposition}[Matched means; C, with E arithmetic]\label{prop:means}
For molecular test weights $f$ and mutant weight $g$, the mean operator of the
augmented process has molecular coordinates
\begin{equation}
 (\mathcal A_\eps(f,g))_i=(Q_ef)_i+b\{(2-\mu_i)(\Lambda f)_i+\mu_ig-f_i\}-d_if_i
 \label{eq:mean}
\end{equation}
and mutant coordinate $(\beta-\delta)g$. It does not involve the joint daughter
law. Consequently $J$ and $I$ have identical type-resolved mean trajectories
under every common deterministic schedule, including the continuation. At
$\eps=0$ the molecular block is $A=Q_e+b(2\Lambda-\Id)-\diag(d)$. Differentiating
\eqref{eq:field} at the all-one vector gives $D\Phi(\one)=\mathcal A_\eps$, with
the \emph{positive} mean-operator sign.
\end{proposition}

\begin{proof}
Differentiate \eqref{eq:offspring} at $(x,y)=(\one,1)$ in the direction $(f,g)$.
Bilinearity and symmetry of $D_D$ with $D_D(\one,\one)=1$ and
$D_D(\one,f)=\Lambda f$ give $(1-\mu_i)\cdot2(\Lambda f)_i$ from the first term,
and $\mu_i(g+(\Lambda f)_i)$ from the second; the sum is
$(2-\mu_i)(\Lambda f)_i+\mu_ig$, which involves $D_D$ only through $\Lambda$.
Adding the transport, death and removal terms of \eqref{eq:field} gives
\eqref{eq:mean}. Both laws induce the same $\Lambda$, so the two mean operators
coincide. Forward mean row vectors obey $m'=m\mathcal A_\eps$ and backward
rewards are columns; the derivative at $\one$ carries a positive sign, and
inserting a negative one corrupts every sensitivity equation below. The exact
rational identity of the two $7\times7$ mean matrices is checked
symbolically \textup{(E)}.
\end{proof}

Proposition~\ref{prop:means} says only that the two laws agree law-to-law for a
fixed common schedule. It does \emph{not} say $AB$ and $BA$ have the same means,
and it does not automatically extend to unconditional means under
population-dependent feedback, where the realised action and the population can
be correlated differently under the two laws.

\subsection{Construction and the probability bridge}

\begin{proposition}[Construction and probability bridge; C]\label{prop:bridge}
For bounded action ranges and a finite initial population the process is
nonexplosive. For deterministic piecewise constant actions, the backward flow of
\eqref{eq:field}--\eqref{eq:fieldM} from a terminal payoff $g\in[0,1]^7$ equals
the actual expected product of $g$ over the terminal population. Chronological
phases $A$ then $B$ give $q=F_A^{10}\bigl(F_B^{10}(g)\bigr)$, where $F^t_\bullet$
is the time-$t$ flow of the corresponding phase field. Fixed founders $z$ give
probability $\prod_iq_i^{z_i}$; a random preparation requires averaging that
product over its actual law.
\end{proposition}

\begin{proof}
Construct bounded per-cell marked Poisson proposals on a countable binary
genealogical tree, with thinning for type changes, death and division. A Yule
process of per-cell rate $b=.1$ dominates the population size. Its finite mean on
compact intervals bounds the expected integrated population and hence the
expected total proposal count, so localisation excludes explosion; this also
covers predictable population feedback, used in Section~\ref{sec:floor}. Under
deterministic rates, future descendant clocks are independent conditional on
daughter types. First-event conditioning gives \eqref{eq:field}: death
contributes the empty product $1$, a type change contributes the new
continuation value, and a division contributes the joint offspring monomial,
averaged over the joint birth law --- this is the unique place where $J$ and $I$
differ. The probability of two events descended from one founder in time $h$ is
$O(h^2)$ uniformly. Bounded convergence and nonexplosion justify the integral
identity and its differentiation. The polynomial field is Lipschitz on the cube,
so uniqueness identifies its solution with the expectation. At a zero coordinate
the field is nonnegative and at a one coordinate it is nonpositive, proving
invariance of the cube. Conditioning at phase boundaries gives chronological
composition. Independent founders multiply only under the stated deterministic
schedule.
\end{proof}

\subsection{The declared continuation and the endpoint}
The objective is eventual total extinction, so the regime after the finite course
is part of the statement. After the deadline $T$ we set $e=.3$, $c=0$ and
$\mu=0$, keeping the $M$ rates fixed.

\begin{proposition}[Clearing continuation and terminal payoff; C, with E arithmetic]
\label{prop:endpoint}
Under this continuation every finite molecular population becomes extinct almost
surely, and the continuation extinction vector is
\begin{equation}
 g=(1,1,1,1,1,1,\rho)=(1,1,1,1,1,1,1/5).\label{eq:terminal}
\end{equation}
Consequently eventual total survival is exactly survival of an $M$ family seeded
by time $T$.
\end{proposition}

\begin{proof}
Let $A_0$ be the mutation-free molecular mean matrix of the clearing source. For
the positive weight $w=(14,11,10,84,43,107)$, exact rational substitution gives
\begin{equation}
 A_0w\le-.09\,w,\qquad -.09\,w-A_0w=(.2,.2,.2,1.56,.78,1.49)\ \ \text{\textup{(E)}}.
 \label{eq:weight}
\end{equation}
By the stopped generator identity and Gronwall, the expected weighted molecular
population decays at rate at least $.09$. Its minimum positive weight is $10$, so
by Markov's inequality the probability of being nonempty at time $t$ is at most
$e^{-.09t}\langle w,z\rangle/10$ and tends to zero; hence almost sure molecular
extinction from any finite configuration. There is no acquisition after $T$, so
no new $M$ appears, and $M$ is autonomous with extinction probability $\rho$ per
initial cell by \eqref{eq:rho}; nonexplosion identifies the limit over
genealogical depths with extinction. Finiteness of the population at $T$ and
conditional independence of the autonomous $M$ families then give
\eqref{eq:terminal} and the stated equivalence.
\end{proof}

Four natural events must be kept apart, since an $M$ cell can appear and later
die. Table~\ref{tab:events} lists them with the terminal payoff computing each,
together with floating diagnostics at the witness parameters which show they are
genuinely different numbers.

\begin{table}[ht]
\centering\small
\caption{Distinct endpoints and their backward constructions. Probabilities are
complementary-law $AB$ diagnostics from one $AA$ founder, evidence \textbf{N};
the certified table is Table~\ref{tab:certificate}.}
\label{tab:events}
\begin{tabular}{llr}\toprule
Event&Backward construction&Diagnostic\\\midrule
Total survival past the deadline $T$&terminal $0$, complement&$.14318$\\
Some $M$ present at $T$&terminal $(\one,0)$, complement&$.03292$\\
Some $M$ ever appeared by $T$&killed molecular field, mutant term deleted&$.03785$\\
Eventual total survival (our objective)&terminal $(\one,\rho)$, complement&$.03042$\\
\bottomrule
\end{tabular}
\end{table}

\section{The dependence mechanism}\label{sec:mechanism}

The certified reversal of Section~\ref{sec:certified} does not rest on this
section: the four enclosed flows prove it outright. But without a mechanism the
result reads as an unexplained numerical coincidence, so we record what actually
drives it. Everything here is a conventional proof, exact algebra, or a finite
exact check.

Throughout this section the terminal molecular payoff is $\one$ and the mutant
coordinate is frozen at its exact fixed point $y\equiv\rho$, which is the
situation of \eqref{eq:terminal}.

\subsection{The first-order response is common}
\begin{proposition}[Rare-acquisition expansion; C]\label{prop:expansion}
On a fixed finite deterministic course, the molecular backward solution admits
the expansion $x_D=\one-\eps h+\eps^2k_D+O(\eps^3)$, in which
\begin{align}
 h'&=Ah+b\kappa(1-\rho),&h(0)&=0,\label{eq:h}\\
 k_D'&=Ak_D+bD_D(h,h)+b\kappa(2-\rho)\Lambda h,&k_D(0)&=0,\label{eq:k}
\end{align}
with $A=Q_e+b(2\Lambda-\Id)-\diag(d)$ and all products with $\kappa$ taken
coordinatewise. The first coefficient $h$ is the same for both daughter laws.
\end{proposition}

\begin{proof}
At $\eps=0$ we have $\mu\equiv0$ and $x\equiv\one$ solves \eqref{eq:field}
exactly, so the expansion starts at $\one$; smooth dependence of a polynomial ODE
on a parameter over a compact finite-horizon domain justifies it. Substituting
$x=\one-\eps h+\eps^2k+O(\eps^3)$ into $D_D$ and using bilinearity, symmetry and
$D_D(\one,f)=\Lambda f$ gives
\[
 D_D(x,x)=1-2\eps\Lambda h+\eps^2\{2\Lambda k+D_D(h,h)\}+O(\eps^3),
\]
while $(\Lambda x)_i=1-\eps(\Lambda h)_i+O(\eps^2)$ and $\mu_i=\eps\kappa_i$.
Collecting the $O(\eps)$ and $O(\eps^2)$ coefficients of \eqref{eq:field} with
$y=\rho$ gives \eqref{eq:h} and \eqref{eq:k}; the detailed collection is in
Appendix~\ref{sec:mechdetail}. The $\eps$-coefficient involves $D_D$ only through
$\Lambda$, which is common, so $h$ is common.
\end{proof}

\subsection{The second order is transported sister covariance}
Both laws are exchangeable with the same marginal, so
$\E_{J,i}[h(Y)]=\E_{J,i}[h(Z)]=(\Lambda h)_i$ and, by \eqref{eq:Ilaw},
\[
 D_{I,i}(h,h)-D_{J,i}(h,h)=(\Lambda h)_i^2-\E_{J,i}[h(Y)h(Z)]
 =-\Cov_{J,i}\bigl(h(Y),h(Z)\bigr).
\]
Subtracting the two copies of \eqref{eq:k} --- the acquisition term
$b\kappa(2-\rho)\Lambda h$ is common and cancels --- gives the transport equation
for $\Delta=k_I-k_J$:
\begin{equation}
 \boxed{\ \Delta'=A\Delta-b\,\Cov_{J}\bigl(h(Y),h(Z)\bigr),\qquad\Delta(0)=0,\ }
 \label{eq:covtransport}
\end{equation}
coordinatewise, and hence, with $U_A(t,s)$ the mean propagator of the
phase-dependent $A$,
\begin{equation}
 \Delta(t)=-\int_0^tU_A(t,s)\,b\,\Cov_{J,s}\bigl(h(Y),h(Z)\bigr)\,ds.
 \label{eq:duhamel}
\end{equation}
Sensitivity states and propagators carry through phase boundaries without reset.
These are finite-horizon expansions; they assert nothing uniformly in time or in
founder number.

\subsection{The correction has a definite sign}
Order the molecular types by
\begin{equation}
 (a,r)\preceq(a',r')\quad:\Longleftrightarrow\quad a\le a'\ \text{and}\ r\ge r',
 \label{eq:order}
\end{equation}
so that ``more activating, fewer repressive'' is higher. This is a partial order
with least element $RR$ and greatest element $AA$, and exactly six covering
pairs: $RR\lessdot UR$, $UR\lessdot UU$, $UR\lessdot AR$, $UU\lessdot AU$,
$AR\lessdot AU$ and $AU\lessdot AA$. Let $\mathcal C$ be the closed convex cone of
vectors that are nonnegative and $\preceq$-increasing.

\begin{lemma}[$h$ stays in the monotone cone; C, with a finite E check]
\label{lem:cone}
At the witness parameters, $h(t)\in\mathcal C$ for every $t\in[0,20]$ and for both
chronological orders.
\end{lemma}

\begin{proof}
The forcing $b\kappa(1-\rho)$ of \eqref{eq:h} is nonnegative, and $\kappa$ is
$\preceq$-increasing: a violation would need $\kappa_i=.4$ and $\kappa_j=.1$ with
$i\preceq j$, hence $i\in\{AU,AA\}$, but every $j$ above $AU=(1,0)$ has $a\ge1$
and $r\le0$ and so lies in $\{AU,AA\}$, and the only $j$ above $AA=(2,0)$ is
$AA$ \textup{(E)}. Since $h(0)=0\in\mathcal C$, it suffices that
$\mathcal C$ is forward invariant for the affine field $f\mapsto Af+b\kappa(1-\rho)$
in each phase. The cone is the polyhedron cut out by $f_{RR}\ge0$ together with
$f_j-f_i\ge0$ over the six covering pairs, so by Nagumo's subtangentiality
criterion invariance holds as soon as each active constraint has nonnegative
outward rate.

For the constraint $f_{RR}\ge0$: if $f\in\mathcal C$ with $f_{RR}=0$ then $f\ge0$
because $RR$ is least, and $A$ is Metzler \textup{(E)}, so
$(Af)_{RR}=A_{RR,RR}\cdot0+\sum_{k\ne RR}A_{RR,k}f_k\ge0$; adding the nonnegative
forcing keeps it nonnegative.

For a covering constraint $i\lessdot j$: if $f\in\mathcal C$ with $f_i=f_j$ we need
$(Af)_j-(Af)_i+b(1-\rho)(\kappa_j-\kappa_i)\ge0$, and the last term is
nonnegative, so it suffices that $(Af)_j\ge(Af)_i$. The slice
$\mathcal C\cap\{f_i=f_j\}$ is again a cone, and the layer-cake decomposition
$f=\int_0^\infty\one_{\{f>t\}}\,dt$ writes any of its elements as a nonnegative
combination of indicators of up-sets, each of which satisfies the same equality
$\one_{\{f>t\}}(i)=\one_{\{f>t\}}(j)$ because $f_i=f_j$. The poset has exactly
seven nonempty up-sets, so the requirement reduces to a finite list of
inequalities. Exact rational evaluation confirms all of them: $32$ inequalities
for phase $A$ and $32$ for phase $B$, all nonnegative \textup{(E)}. Both phases
therefore leave $\mathcal C$ invariant, and the two phases compose because the cone
is the same. The explicit inequalities are in Appendix~\ref{sec:mechdetail}.
\end{proof}

\begin{theorem}[The independent approximation is optimistic at leading order; C]
\label{thm:sign}
Under the hypothesis of Lemma~\ref{lem:cone},
$\Cov_{J,i}(h(Y),h(Z))\le0$ for every molecular type $i$ and every time, and
therefore $\Delta=k_I-k_J\ge0$ coordinatewise. Ignoring sister dependence
over-estimates the eventual extinction probability at leading order in $\eps$,
equivalently it under-estimates the probability of rescue.
\end{theorem}

\begin{proof}
Fix a mother of type $(a,r)$ and let $\xi_1,\dots,\xi_a$ and
$\zeta_1,\dots,\zeta_r$ be the independent fair allocation bits of
\eqref{eq:Jlaw}, so that $Y=(\sum_m\xi_m,\sum_m\zeta_m)$ and
$Z=(a-\sum_m\xi_m,\,r-\sum_m\zeta_m)$. Introduce the independent bit vector
\[
 W=(\xi_1,\dots,\xi_a,\,1-\zeta_1,\dots,1-\zeta_r)\in\{0,1\}^{a+r},
\]
whose coordinates are again independent and fair. Increasing a coordinate of $W$
either increases the activating count of $Y$ or decreases its repressive count,
so $Y$ is $\preceq$-increasing in $W$; simultaneously it decreases the activating
count of $Z$ or increases its repressive count, so $Z$ is $\preceq$-decreasing in
$W$. By Lemma~\ref{lem:cone}, $h$ is $\preceq$-increasing, so $h(Y)$ is a
coordinatewise increasing and $h(Z)$ a coordinatewise decreasing function of the
independent vector $W$. Harris' correlation inequality for product measures
\cite{harris1960} gives $\Cov(h(Y),h(Z))\le0$.

The forcing $-b\Cov_J(h(Y),h(Z))$ in \eqref{eq:covtransport} is therefore
nonnegative. Each phase matrix $A$ is Metzler, so its propagator $U_A(t,s)$ is
entrywise nonnegative, and \eqref{eq:duhamel} with $\Delta(0)=0$ gives
$\Delta\ge0$. Since $x_D=\one-\eps h+\eps^2k_D+O(\eps^3)$ with $h$ common, the
first term of $x_I-x_J$ that does not cancel is $\eps^2\Delta$, and it is
nonnegative; by Propositions~\ref{prop:bridge} and \ref{prop:endpoint} $x$ is
the eventual extinction probability.
\end{proof}

Two qualifications matter. First, Theorem~\ref{thm:sign} is a statement about the
fixed source above, with constant division rate and the stated monotone death and
acquisition profiles; state-dependent division would break Lemma~\ref{lem:cone},
whose finite check is specific to these two phase matrices. Second, and
decisively for our purpose, a nonnegative dependence error \emph{for each
schedule separately} does not decide which schedule has the larger error. The
reversal depends on that difference, which the next two results address.

\subsection{An exact finite-amplitude response identity}
The expansion above explains where the effect comes from but cannot certify it.
At the witness value $\eps=1/10$ the best remainder bound we obtained from a
generic third-derivative estimate exceeds $600$, seven orders of magnitude larger
than the $10^{-6}$ decision (Appendix~\ref{sec:mechdetail}). That is a failed
sufficient bound, not a false expansion. The following identity replaces it and
is exact at every $\eps$.

\begin{proposition}[Exact dependence response; C]\label{prop:exact}
For a common terminal molecular payoff let $z=x_I-x_J$ and put
\begin{align}
 B(t)f&=Q_ef-(d+b)f+b(1-\mu)D_I(x_I+x_J,f)+b\mu\rho\Lambda f,\label{eq:Bop}\\
 C(t)&=b(1-\mu)\bigl\{D_I(x_J,x_J)-D_J(x_J,x_J)\bigr\}
      =-b(1-\mu)\Cov_J\bigl(x_J(Y),x_J(Z)\bigr).\label{eq:Cforce}
\end{align}
Then $z'=B(t)z+C(t)$ with $z(0)=0$, so
$z(t)=\int_0^tU_B(t,s)C(s)\,ds$. The operator $B(t)$ is Metzler on the invariant
cube, so $U_B$ is nonnegative, and the response to $|C|$ bounds the absolute
dependence error.
\end{proposition}

\begin{proof}
Subtract the two copies of \eqref{eq:field} and add and subtract
$D_I(x_J,x_J)$. Bilinearity and symmetry of $D_I$ give
$D_I(x_I,x_I)-D_I(x_J,x_J)=D_I(x_I+x_J,x_I-x_J)$, which is the third term of
\eqref{eq:Bop} applied to $z$; the transport, death, removal and mutant terms are
linear in $x$ and common to both laws, giving the remaining terms of
\eqref{eq:Bop}. What is left over is exactly \eqref{eq:Cforce}, and the second
form of $C$ follows as in \eqref{eq:covtransport}. All coefficients of $Q_e$,
$D_I$ and $\Lambda$ off the diagonal are nonnegative and $x_I+x_J\ge0$ on the
cube, so $B(t)$ is Metzler.
\end{proof}

Proposition~\ref{prop:exact} isolates the whole dependence error: the common
acquisition terms enter only through the linear response $B$, and cancel from the
forcing $C$. Note that \eqref{eq:Cforce} is again a sister covariance, now of the
full finite-$\eps$ solution rather than of the first-order coefficient.

\subsection{From the mechanism to the certified reversal}
Write $E_s=q_I(s)-q_J(s)$ for the dependence error of schedule $s$ and
$G_D=q_D(AB)-q_D(BA)$ for the signed order gap of law $D$. Then, trivially but
usefully,
\begin{equation}
 G_J=G_I-(E_{AB}-E_{BA}).\label{eq:gapidentity}
\end{equation}
This is the bridge between mechanism and certificate. Theorem~\ref{thm:sign} says
both $E_{AB}$ and $E_{BA}$ are positive at leading order, and
Proposition~\ref{prop:exact} says each is a transported covariance; because the
covariance is generated during one pulse and then propagated through the other,
the two schedules transport it differently, and their difference is what
\eqref{eq:gapidentity} feeds into the order gap. Section~\ref{sec:certified}
certifies that this difference exceeds the order margin, which is exactly the
condition for a reversal. Figure~\ref{fig:decision}(c) displays the leading-order
transport for the two schedules; at $\eps=1/10$ it reproduces each certified
error, and their difference, to within about $18\%$, which is the expected size
of the neglected $O(\eps^3)$ term \textup{(N)}.

Neither the covariance sign alone nor any qualitative asymptotic argument proves
the finite witness. The enclosed flows do.

\section{A certified equal-exposure decision reversal}\label{sec:certified}

\subsection{Actions and exposure accounting}
Let $\eps=1/10$, so the acquisition probability per division is $.04$ for the
protected types and $.01$ otherwise. Fix the two actions
\[
 A:(e,c)=(3/10,\,0),\qquad B:(e,c)=(1/100,\ 120553/10^6),
\]
each applied for ten time units and followed by the continuation of
Proposition~\ref{prop:endpoint}. Action $A$ erases marks aggressively without
extra killing; action $B$ kills without erasing. Both orders have deadline
$T=20$, intrinsic-erasure increment exposure $2.9$ above the baseline $.01$, and
additional molecular-death exposure $1.20553$. The $M$ rates are unaffected by
either action.

These are equal integrals of model rate increments. They are not equal measured
drug masses, equal toxicity burdens, equal pharmacokinetic exposures or equal
healthy-cell constraints; the equality is an assumption of the comparison, chosen
so that no difference in the total amount of intervention can explain the result.

\subsection{The counterexample}

\begin{theorem}[Matched means, opposite decisions; C with E arithmetic]
\label{thm:decision}
For every positive integer $n$ of identical $AA$ founders, complementary
inheritance $J$ gives strictly larger eventual total-extinction probability under
$BA$ than under $AB$, while independent-marginal inheritance $I$ gives strictly
larger probability under $AB$ than under $BA$. At the central parameters each
one-founder order margin exceeds $1.6\times10^{-6}$ in magnitude. Both
preferences persist for every
\begin{equation}
 .120545\le c_B\le.120561,\label{eq:neighborhood}
\end{equation}
all other quantities held fixed.
\end{theorem}

\begin{proof}
The exact outward arithmetic of Appendix~\ref{sec:certificate} places the four
one-founder extinction probabilities in the rational intervals of
Table~\ref{tab:certificate}, whence
\begin{align}
 -1.702\times10^{-6}&<q_J(AB)-q_J(BA)<-1.701\times10^{-6},\label{eq:gapJ}\\
  1.770\times10^{-6}&<q_I(AB)-q_I(BA)< 1.771\times10^{-6}.\label{eq:gapI}
\end{align}
Propositions~\ref{prop:bridge} and \ref{prop:endpoint} identify these $AA$
coordinates with eventual total extinction from one $AA$ founder. All four
numbers lie strictly in $(0,1)$, so raising them to the power $n$ preserves the
strict orders. Proposition~\ref{prop:means} gives equality of all mean
trajectories between the two laws, schedule by schedule. The interval assertion
is Lemma~\ref{lem:curvature} below.
\end{proof}

\begin{table}[ht]
\centering
\caption{Rational outward intervals for eventual total extinction from one $AA$
founder, at the central parameters. The exact dyadic intervals in the
reproduction files are tighter; each order-gap enclosure has width below
$4\times10^{-16}$.}
\label{tab:certificate}
\begin{tabular}{lll}\toprule
Law and order&Lower endpoint&Upper endpoint\\\midrule
$J$, $AB$&.969579866104740&.969579866104742\\
$J$, $BA$&.969581567445558&.969581567445561\\
$I$, $AB$&.969623237350068&.969623237350071\\
$I$, $BA$&.969621466763777&.969621466763781\\\bottomrule
\end{tabular}
\end{table}

\subsection{Curvature turns four endpoint flows into a continuum}
The interval statement cannot be obtained by sampling. We prove it from two
enclosed endpoint gaps, the preserved central certificate and an analytic bound
on the second derivative of the gap in the parameter. We write the argument out
in full, because the intermediate step --- how a secant becomes a uniform
derivative bound --- is what makes the conclusion valid off the sampled points.

\begin{lemma}[Curvature and the certified interval; C with E arithmetic]
\label{lem:curvature}
Let $G_D(c)=q_D(AB;c)-q_D(BA;c)$ as a function of the additional death rate $c$
in pulse $B$, all else fixed. Then $|G_D''|\le6912$ uniformly on the nonnegative
parameter range, and $G_J<0<G_I$ throughout \eqref{eq:neighborhood}.
\end{lemma}

\begin{proof}
Work with the six molecular coordinates, the mutant coordinate being frozen at
$\rho$. Write the molecular backward field along the schedule as $f(t,x,c)$, let
$\chi_B(t)$ indicate the phase carrying the varied parameter, and set $J_x=D_xf$,
$u=\partial_cx$, $v=\partial_c^2x$. Since $c$ enters \eqref{eq:field} only through
$d_i=d_i^0+c$ in the term $d_i(1-x_i)$, we have $\partial_cf=\chi_B(\one-x)$,
$\partial_c\partial_xf=-\chi_B\Id$ and $\partial_c^2f=0$. Differentiating the
flow equation once and twice in $c$ gives
\begin{align}
 u'&=J_xu+\chi_B(t)(\one-x),&u(0)&=0,\label{eq:u}\\
 v'&=J_xv+D_{xx}f[u,u]-2\chi_B(t)u,&v(0)&=0.\label{eq:v}
\end{align}
The factor two in \eqref{eq:v} arises because
$\partial_cJ_x=D_{xx}f[u,\cdot]-\chi_B\Id$ contributes $-\chi_Bu$ in addition to
the $-\chi_Bu$ obtained by differentiating the forcing of \eqref{eq:u}.
Sensitivities pass through the phase boundary without resetting, and pulse
durations are fixed, so there is no moving-boundary derivative.

Three uniform bounds on the cube are needed, all by direct inspection of
\eqref{eq:source}--\eqref{eq:offspring} \textup{(E)}. The Jacobian $J_x$ is
Metzler with molecular row sums at most $b-d_i-b\mu_i(2-\rho)\le b-d_i\le.09$, so
its logarithmic sup norm is at most $.09$. The state Hessian satisfies
$\|D_{xx}f[u,u]\|_\infty\le2b\|u\|_\infty^2=.2\|u\|_\infty^2$, since the only
quadratic term is $b(1-\mu_i)D_i(x,x)$ and $D_i$ has unit total mass. At the
molecular vector $\one$ with mutant coordinate $\rho$ the field equals
$-b\mu_i(1-\rho)$, of sup norm at most $b\eps\kappa_{\max}(1-\rho)=.0032$; this
is \emph{not} evaluated at the full seven-dimensional all-one vector, where the
field vanishes.

Variation of constants applied to $\one-x$, then to \eqref{eq:u}, gives for
$0\le t\le20$
\begin{equation}
 \|\one-x(t)\|_\infty\le.0032\,te^{.09t},\qquad
 \|u(t)\|_\infty\le\int_0^te^{.09(t-s)}.0032\,se^{.09s}ds=.0016\,t^2e^{.09t}<6,
 \label{eq:ubound}
\end{equation}
the last step using $e^{1.8}<9$, which follows from $e<3$ and hence
$e^{1.8}<e^2<9$. Feeding \eqref{eq:ubound} into \eqref{eq:v},
\begin{equation}
 \|v(20)\|_\infty\le\int_0^{20}e^{.09(20-s)}\bigl\{.2\|u(s)\|_\infty^2
   +2\|u(s)\|_\infty\bigr\}ds\le20\cdot9\{.2\cdot36+2\cdot6\}=3456.
 \label{eq:vbound}
\end{equation}
Each $G_D$ is a difference of two such schedule coordinates, so
$|G_D''|\le2\cdot3456=6912$ uniformly, and the bound holds on both phases and on
the whole nonnegative parameter range.

Now the endpoints. Exact arithmetic on four additional flows with identical
source equations gives
\begin{equation}
 -5.114\times10^{-7}<G_J(.120545)<-5.113\times10^{-7},\qquad
  5.830\times10^{-7}<G_I(.120561)<5.831\times10^{-7}\ \ \text{\textup{(E)}}.
 \label{eq:endpoints}
\end{equation}
For law $J$ form the secant of $G_J$ between $.120545$ and the central value
$.120553$, a span of $8\times10^{-6}$; combining \eqref{eq:endpoints} with
\eqref{eq:gapJ} encloses that secant slope in
$(-.1487475,-.1487474)$. By the mean value theorem the secant equals
$G_J'(\xi)$ at some $\xi$ strictly inside that span, hence inside
\eqref{eq:neighborhood}. Every point $c$ of \eqref{eq:neighborhood} satisfies
$|c-\xi|\le1.6\times10^{-5}$, so the curvature bound gives
$|G_J'(c)-G_J'(\xi)|\le6912\cdot1.6\times10^{-5}=.110592$ and therefore
\[
 G_J'(c)<-.1487474+.110592=-.03815\qquad\text{for all $c$ in
 \eqref{eq:neighborhood}}.
\]
The same computation for $I$, using the secant between $.120553$ and $.120561$
together with \eqref{eq:gapI}, gives $G_I'(c)<-.03785$ on the same interval.
Both gaps are therefore strictly decreasing across \eqref{eq:neighborhood}. Since
$G_J$ is negative at the left endpoint and decreasing, it is negative throughout;
since $G_I$ is positive at the right endpoint and decreasing, it is positive
throughout. No sampling and no root bracket are used.
\end{proof}

\begin{remark}
The interval \eqref{eq:neighborhood} is an \emph{inner} certified interval, not
the maximal reversal set: it is bounded by where we placed endpoint flows. Its
width $1.6\times10^{-5}$ is $8000$ times the width of the elementary Lipschitz
neighbourhood ($2\times10^{-9}$) available before this argument, which measures
an improvement in the certificate and not an amplification of any biological
effect. The margins at the endpoints are correspondingly smaller than the central
margins. Determining the maximal set would require either many more endpoint
flows or a validated root enclosure; neither is needed here.
\end{remark}

\subsection{The approximation error against the decision margin}
Ignoring sister dependence changes the extinction coordinate by
\begin{equation}
 E_{AB}\approx4.3371245\times10^{-5},\qquad
 E_{BA}\approx3.9899318\times10^{-5},\qquad
 E_{AB}-E_{BA}\approx3.4719271\times10^{-6},
 \label{eq:errors}
\end{equation}
all three enclosed exactly \textup{(E)}. Both absolute errors are small ---
about four parts in $10^{5}$ of a probability close to $.97$ --- and they have
the sign predicted by Theorem~\ref{thm:sign}. But their \emph{difference} exceeds
the approximate order margin $1.771\times10^{-6}$, and by
\eqref{eq:gapidentity} that is precisely what flips the sign of the gap.
Absolute probability accuracy and decision accuracy are therefore different
things. This result isolates decision vulnerability near a tie; it does not claim
a large population-level benefit.

\begin{figure}[t]
\centering\includegraphics[width=\linewidth]{figures/decision.pdf}
\caption{(a) The shaded interval is certified by two endpoint enclosures, the
preserved central certificate and the analytic curvature bound of
Lemma~\ref{lem:curvature}; the smooth curves are floating-point diagnostics and
their zero crossings, near $.1205416$ and $.1205649$, are neither certified root
brackets nor a claim about a maximal reversal region. (b) Curves use midpoints of
the enclosed one-founder probabilities; the unique maxima at $32$ founders are
established by exact rational increment signs, not by the plotted samples.
(c) Leading-order dependence error $\eps^2\Delta_{AA}$ transported backward from
the deadline, from \eqref{eq:covtransport}, with the phase boundary at $s=10$;
stars mark the certified finite-$\eps$ values $E_s$ of \eqref{eq:errors}. The
curves cross: the schedules differ in how the covariance forcing generated in one
pulse is propagated through the other.}
\label{fig:decision}
\end{figure}

\section{What the approximation costs a decision maker}\label{sec:decision}

\subsection{The direct misspecification cost}
The cleanest operational quantity involves no hypothetical randomisation over
models. Suppose the complementary law $J$ is the true source and a modeller,
having validated the means, adopts the independent approximation and follows its
recommendation $AB$. The loss relative to the recommendation $BA$ that $J$ would
have given is
\begin{equation}
 R_{I\to J}=q_J(BA)-q_J(AB)=m_J\in
 \bigl(1.70134081854\times10^{-6},\,1.70134081866\times10^{-6}\bigr)\ \ \text{\textup{(E)}}.
 \label{eq:direct}
\end{equation}
This is simultaneously lost extinction probability and added eventual-survival
probability. We write $m_J=q_J(BA)-q_J(AB)>0$ and $m_I=q_I(AB)-q_I(BA)>0$ for the
two order margins throughout.

\subsection{When a schedule ranking is protected}
\begin{corollary}[A replacement for unqualified mean-based decisions; C]
\label{cor:criterion}
If two schedules admit valid bounds $|q_J(s)-q_I(s)|\le E_s$ and
$q_I(A)-q_I(B)>E_A+E_B$, then $q_J(A)>q_J(B)$. Signed bounds on $E_A-E_B$, when
available, are sharper and may be used instead. Conversely, no decision rule
using only the complete type-resolved mean trajectories can always select the
better of our two schedules for both inheritance laws.
\end{corollary}

\begin{proof}
The first assertion is the triangle inequality. For the second, the mean data are
identical for the two models by Proposition~\ref{prop:means}, so a rule using only
those data returns the same answer in both; Theorem~\ref{thm:decision} requires
opposite answers.
\end{proof}

The worked instance is instructive. The independent-model central margin lies in
$(1.770,1.771)\times10^{-6}$ and the paired correction $E_{AB}-E_{BA}$ in
$(3.471,3.473)\times10^{-6}$, so subtracting the signed correction certifies the
complementary preference for $BA$. Had only a symmetric bound of
$\pm4\times10^{-6}$ on that correction been available, the corrected margin
interval would contain zero and the defensible output would be ``order
unresolved at this approximation-error level'' \textup{(E)}. This concerns model
approximation error and is not an empirical confidence interval; empirical
uncertainty would have to be modelled separately.

For a fixed heterogeneous preparation, telescoping the product of coordinates
gives $|P_J-P_I|\le\sum_iz_i|q_{J,i}-q_{I,i}|$. An ordering for identical $AA$
founders proves no ordering for an arbitrary mixture.

\subsection{Founder scaling}
\begin{proposition}[Unique founder maximum; C with E arithmetic]\label{prop:founders}
Fix a law and let $0<l<h<1$ be its worse and better one-founder extinction
probabilities. The advantage $A_n=h^n-l^n$ is positive, tends to zero, and
satisfies $A_n\le n(h-l)h^{n-1}$. Its increment $A_{n+1}-A_n=l^n(1-l)-h^n(1-h)$
changes sign at most once. For both laws the increment is positive at $n=31$ and
negative at $n=32$, so both advantages attain their unique integer maximum at
$32$ identical $AA$ founders, with values in
$(2.08950953,2.08950954)\times10^{-5}$ for $J$ and
$(2.17745028,2.17745029)\times10^{-5}$ for $I$.
\end{proposition}

\begin{proof}
Positivity and the factorisation $A_n=(h-l)\sum_{j<n}h^{n-1-j}l^j\le n(h-l)h^{n-1}$
are immediate, as is $A_n\to0$. Dividing the increment by $h^n$ gives
$(l/h)^n(1-l)-(1-h)$, and $(l/h)^n$ strictly decreases, so the increment changes
sign at most once. The two increment signs are evaluated in exact rational
arithmetic from the endpoints of Table~\ref{tab:certificate} \textup{(E)}, which
avoids selecting an integer by rounded logarithms.
\end{proof}

A larger founder count therefore does not amplify the effect indefinitely; it
changes achieved risk without changing the preferred schedule
(Figure~\ref{fig:decision}(b)).

\subsection{A minimax reading of the information limitation}
The next proposition is a statement about a class of models that the permitted
information cannot separate. It is secondary to \eqref{eq:direct} and we state it
as such.

\begin{proposition}[Regret under mean-indistinguishable information; C with E arithmetic]
\label{prop:regret}
A randomised decision rule that sees only information common to the two models,
including their complete type-resolved mean trajectories, has worst-case expected
extinction regret at least
\[
 R_*=\frac{m_Jm_I}{m_J+m_I}\in
 \bigl(8.67636512\times10^{-7},\,8.67636514\times10^{-7}\bigr),
\]
and the bound is attained in this two-model problem by choosing $AB$ with
probability $m_I/(m_J+m_I)$.
\end{proposition}

\begin{proof}
Such a rule selects $AB$ with the same probability $p$ under both laws. Its
regrets are $pm_J$ and $(1-p)m_I$. If $p\ge m_I/(m_J+m_I)$ the first is at least
$R_*$, otherwise the second is; the displayed probability equalises them. The
function $(s,t)\mapsto st/(s+t)$ increases in each positive argument, so
substituting the endpoints of the certified intervals for $m_J$ and $m_I$ gives
the enclosure \textup{(E)}. The randomisation selects an entire deterministic
course independently of the population; the claim does not restrict a rule that
observes sister pairs or any other information distinguishing the two laws.
\end{proof}

\begin{remark}
The premise is an information restriction, not a physical randomisation of
nature. If the conservative law $J$ is already known to hold, the two-world
premise does not apply and the relevant loss is the deterministic
\eqref{eq:direct}, which is about twice $R_*$. Halving the number by
randomisation does not reduce the cost of the deterministic mistake under $J$.
Nor is $p_*$ a dosing recommendation. Note also that an optimiser working inside
model $I$ uses that model's assumed offspring law as well as the mean data; the
actual modelling failure is treating validated means as sufficient warrant for
that additional assumption.
\end{remark}

\section{A survival floor for every admitted feedback course}\label{sec:floor}

The reversal compares two fixed schedules. It is natural to ask what any
controller could achieve, including one that adapts to the observed population.
The next theorem answers this in the direction that matters --- it is a
feasibility obstruction --- and it holds without assuming independence between
founder outcomes, which fails under feedback.

\begin{theorem}[Scalar feedback envelope; C with E arithmetic]\label{thm:floor}
On $[0,T]$ let actions be bounded and predictable, let division have constant
rate $b>0$, and let $d_i\le d_{\max}$ and $\mu_i\ge\mu_{\min}>0$. Retain the
normalised binary daughter law and the one-mutant-daughter convention. Let $M$
have autonomous rates $\beta>\delta>0$ with $\rho=\delta/\beta$, and assume
$d_{\max}\ge b\rho$. Let the post-course continuation create no new $M$, clear
every finite molecular population almost surely, and retain the autonomous $M$
rates. Define $q$ by
\begin{equation}
 q'=d_{\max}(1-q)+b\bigl\{(1-\mu_{\min})q^2+\mu_{\min}\rho q-q\bigr\},
 \qquad q(0)=1.\label{eq:scalar}
\end{equation}
Then from $n$ molecular founders of arbitrary types and $m$ initial mutants every
such feedback course satisfies
\begin{equation}
 \Prb(\text{eventual total survival})\ge1-q(T)^n\rho^m.\label{eq:floor}
\end{equation}
\end{theorem}

\begin{proof}
The scalar field of \eqref{eq:scalar} is $-b\mu_{\min}(1-\rho)<0$ at $q=1$ and
$(1-\rho)(d_{\max}-b\rho)\ge0$ at $q=\rho$, so $[\rho,1]$ is invariant and
$q\in[\rho,1]$ throughout. Evaluate the actual molecular field at the constant
molecular vector with all coordinates equal to $q$ and mutant coordinate $\rho$.
All molecular jumps cancel because $Q_e$ is conservative, and every normalised
two-daughter law satisfies $D(q\one,q\one)=q^2$ --- this is where the comparison
becomes law-independent. The difference between the scalar field and the actual
one is
\[
 (d_{\max}-d_i)(1-q)+b(\mu_i-\mu_{\min})q(q-\rho)\ge0.
\]
For a population $z$ set $V(t,z)=q(T-t)^{\sum_{i\in S}z_i}\rho^{z_M}$. Direct
evaluation of each jump gives, for the currently chosen predictable action $u$,
\[
 (\partial_t+\mathcal G_u)V=V\sum_i\frac{z_i}{q_i}\{\Phi_i^u(q)-q_i'\}\le0,
 \qquad q_M=\rho.
\]
This is a conditional generator identity for a population test function and
requires no independence between founder outcomes. Stop at bounded population and
proposal count, apply the stopped Dynkin inequality, and remove the stopping by
nonexplosion (Proposition~\ref{prop:bridge}) and bounded convergence, using
$0\le V\le1$. At $T$, $V(T,Z_T)=\rho^{Z_M(T)}$ is the conditional probability of
eventual extinction under the declared continuation, since that continuation
clears the molecular population and creates no new $M$. Taking expectations gives
\eqref{eq:floor}.
\end{proof}

On the admitted action rectangle $.01\le e\le.3$, $0\le c\le.121$ with $\eps=.1$,
we have $b=.1$, $d_{\max}=.421$, $\mu_{\min}=.01$ and $\rho=.2$, and
$d_{\max}\ge b\rho$ holds comfortably. Writing $a=1-q$ reduces \eqref{eq:scalar}
exactly to the Riccati equation
\begin{equation}
 a'=.0008-.3228\,a-.099\,a^2,\qquad a(0)=0\ \ \text{\textup{(E)}},
 \label{eq:riccati}
\end{equation}
and an exact scalar enclosure gives
\begin{equation}
 .0024725782<a(20)<.0024725784\ \ \text{\textup{(E)}}.\label{eq:concretefloor}
\end{equation}
A one-founder eventual-survival target below $.00247$ is therefore impossible in
this actuator class, whatever the feedback policy. The risk actually achieved by
complementary $BA$ is about $.03041843$. The gap between the achieved risk and
the universal necessary bound is a factor of about twelve: the floor is a genuine
obstruction, but neither global optimality nor near-optimality of the exhibited
schedule follows from it. The floor pertains to the stated endpoint and action
class; an arbitrary new continuation is not automatically covered.

\section{How much continuation perturbation the decision tolerates}
\label{sec:continuation}

The baseline clearing continuation runs indefinitely and is not drug-free:
maintaining $e=.3$ above the baseline $.01$ has infinite integrated erasure
increment if the continuation is charged to the exposure budget. The equal
exposures of Section~\ref{sec:certified} compare the two finite courses. The two
propositions below quantify two separate departures from the idealised
continuation, both at the central course parameters. They are separate
scenarios; combining them, altering the $M$ dynamics, or adding population
interactions requires a further argument.

Both bounds are one-sided, which is what makes them efficient. Writing
$\widehat q_s$ for the perturbed eventual extinction probability, in each case
$0\le q_s-\widehat q_s\le\theta_s$, so only $\theta_{BA}$ can erode the $J$
preference $m_J$ and only $\theta_{AB}$ can erode the $I$ preference $m_I$.

\subsection{Residual acquisition after the course}
\begin{proposition}[Residual acquisition; C with E arithmetic]\label{prop:residual}
Keep $e=.3$, $c=0$ after $T$, but replace zero acquisition by
$\mu_i=\eta\kappa_i$. For either inheritance law and either course $s$,
\[
 0\le q_s(0)-q_s(\eta)\le\eta U_s,\qquad
 U_{AB}<.050002,\quad U_{BA}<.053172.
\]
Both central order preferences persist for every $0\le\eta\le3\times10^{-5}$.
\end{proposition}

\begin{proof}
Let $A_0$ be the mutation-free molecular clearing mean matrix and $f=b\kappa$.
Matrices act on column rewards, forward mean rows obeying $m'=mA$. The perturbed
molecular mean matrix is $A_\eta=A_0-\diag(b\eta\kappa)\Lambda\le A_0$ entrywise
and is Metzler, so semigroup comparison together with \eqref{eq:weight} preserves
almost sure molecular clearing and bounds the molecular occupation means. Define
the nonnegative column
\[
 v=\int_0^\infty e^{tA_0}f\,dt=-A_0^{-1}f,
\]
solved exactly in rationals \textup{(E)}. Integrating the acquisition intensity
against the dominated molecular mean and conditioning on each seed's autonomous
descendants, the expected number of additional surviving resistant families from
molecular type $i$ is at most $(1-\rho)\eta v_i$; a union bound converts that
expectation into a probability bound. Coupling with zero acquisition shows
extinction can only decrease, because baseline molecular descendants already
clear and the additional surviving mutant families are the only new source of
survival. Conditioning at the deadline and writing $A_A,A_B$ for the
finite-course molecular mean matrices gives
\[
 U_{AB}=(1-\rho)\bigl[e^{10A_A}e^{10A_B}v\bigr]_{AA},\qquad
 U_{BA}=(1-\rho)\bigl[e^{10A_B}e^{10A_A}v\bigr]_{AA},
\]
enclosed by validated linear flows as described in
Appendix~\ref{sec:linear}. Since the corrections are one-sided, it suffices that
$\eta U_{BA}<m_J$ and $\eta U_{AB}<m_I$. At $\eta=3\times10^{-5}$ these read
$1.59515\times10^{-6}<1.70134\times10^{-6}$ and
$1.50006\times10^{-6}<1.77058\times10^{-6}$ \textup{(E)}.
\end{proof}

The binding constraint is the first, and the same bounds fail at
$\eta=3.2\times10^{-5}$, so $3\times10^{-5}$ is close to what this argument can
give. It allows post-course acquisition probabilities up to $1.2\times10^{-5}$
per division, still far below the in-course value $.04$; the improvement is a
sharper sufficient radius, not a demonstration of broad biological robustness.
Under residual acquisition the molecular continuation extinction coordinates fall
below one and become law dependent, so \eqref{eq:terminal} is not reused as their
exact value.

\subsection{Finite clearing and its exposure cost}
\begin{proposition}[Finite clearing; C with E arithmetic]\label{prop:clearing}
Apply the baseline mutation-free clearing source for only $H$ time units after
$T$, and then a common specified source that retains autonomous $M$, has no
immigration, and differs from baseline only through the remaining molecular
cells. If $\widetilde q_s$ is the resulting eventual extinction probability, then
\[
 0\le q_s-\widetilde q_s\le C_se^{-.09H},\qquad
 C_{AB}<1.430424,\quad C_{BA}<1.452994,
\]
and $H=152$ preserves both central preferences. The additional common erasure
increment exposure is $.29H=44.08$.
\end{proposition}

\begin{proof}
The contraction weight \eqref{eq:weight} bounds the probability that any
molecular cell remains at $T+H$ by $C_se^{-.09H}$, with
$C_{AB}=[e^{10A_A}e^{10A_B}w]_{AA}/10$ and similarly for $BA$; the constants are
enclosed by validated linear flows (Appendix~\ref{sec:linear}).

For the sign, condition on the configuration at $T+H$. Under the baseline
continuation the conditional eventual extinction probability is exactly
$\rho^{Z_M}$, by Proposition~\ref{prop:endpoint}. Under the alternate source,
total extinction implies extinction of the autonomous descendants of the $M$
cells already present, an event of probability exactly $\rho^{Z_M}$ and
unaffected by any later molecular behaviour; hence the conditional probability is
at most $\rho^{Z_M}$. Therefore $\widetilde q_s\le q_s$. On the event that no
molecular cell remains at $T+H$ the two continuations couple exactly, since the
alternate source differs from baseline only through remaining molecular cells and
admits no immigration, so the difference is supported on the complementary event
and is bounded by its probability.

Because the error is one-sided, it suffices that $C_{BA}e^{-.09H}<m_J$ and
$C_{AB}e^{-.09H}<m_I$. Bounding $e^{.09H}$ from below by a truncation of its
nonnegative power series, at $H=152$ these read
$1.66386\times10^{-6}<1.70134\times10^{-6}$ and
$1.63801\times10^{-6}<1.77058\times10^{-6}$ \textup{(E)}. The same bounds fail at
$H=151$.
\end{proof}

This sign argument requires that later molecular behaviour cannot alter the
autonomous descendants of the $M$ cells already present. It must not be extended
to a continuation violating that assumption. Under the cruder two-sided estimate
$|\widetilde q_s-q_s|\le C_se^{-.09H}$ with the errors summed, the same margins
require $H=162$; the one-sided argument saves ten time units.

The honest reading is twofold, and both halves should be stated together. Finite
clearing \emph{is} realisable: the counterexample does not secretly depend on an
infinite intervention. But the sufficient duration is expensive. The added common
erasure increment exposure $44.08$ is exactly $15.2$ times the finite-course
increment $2.9$, and the clearing phase lasts $7.6$ times as long as the course
itself. The initial sequencing still changes how many resistant families are
seeded, and clearing does not erase that difference because $M$ is unaffected by
it; but the common burden substantially weakens the practical attractiveness of
the construction. No optimality of $H=152$ is asserted, and $H$ is a sufficient
duration derived from a bound, not a physical minimum. Replacing the global
weight estimate by an enclosed expected molecular cell count at $T+H$ would
tighten it; no such optimisation is needed for the methodological theorem.

\section{Dimensionless parameters, data requirements and limits}\label{sec:units}

Put $\tau=bt$ with reference division rate $b=.1$. Table~\ref{tab:units} gives
every source coefficient in these units. Each pulse then has dimensionless
duration one and the course has $\tau=2$.

\begin{table}[ht]
\centering\small
\caption{Complete dimensionless synthetic specification. Rates divide by $b$;
probabilities and mark counts do not. Erasure $e$ and additional death $c$ are
model rates, not drug concentrations.}\label{tab:units}
\begin{tabular}{lll}\toprule
Quantity&Dimensional value or formula&Dimensionless value\\\midrule
Activating write rate&$u(.01+a/2)$&$u(.1+5a)$\\
Repressive write rate&$u(.01+r/2)$&$u(.1+5r)$\\
Activating erase rate&$a(e+r/4)$&$a(e/b+2.5r)$\\
Repressive erase rate&$r(e+a/4)$&$r(e/b+2.5a)$\\
Molecular division&$.1$&$1$\\
Base death, $AU,AA$&$.01$&$.1$\\
Base death, other types&$.3$&$3$\\
Pulse $A$: $(e,c)$&$(.3,0)$&$(3,0)$\\
Pulse $B$: $(e,c)$&$(.01,.120553)$&$(.1,1.20553)$\\
Pulse duration; deadline&$10;20$&$1;2$\\
Clearing: $(e,c,\mu)$&$(.3,0,0)$&$(3,0,0)$\\
Mutant birth; death&$.1;.02$&$1;.2$\\
Course acquisition $AU,AA$&$\eps\kappa=.1(.4)$&$.04$ per division\\
Course acquisition, other&$\eps\kappa=.1(.1)$&$.01$ per division\\
Partition probability&$1/2$ per old mark&$1/2$\\
Feedback erasure range&$[.01,.3]$&$[.1,3]$\\
Feedback additional death&$[0,.121]$&$[0,1.21]$\\\bottomrule
\end{tabular}
\end{table}

The correct reading of these numbers is narrow. The initial two-pulse course is
short relative to the reference division timescale: its integrated division
hazard is two. Calendar duration, cell-cycle statistics, acquisition mechanism
and intervention response all remain uncalibrated. It would be wrong to call the
course ``two cell generations'': division is an exponential clock rather than a
synchronised cycle, so $bT=2$ is the expected number of divisions along an
immortal tracked lineage, while death and conditioning on survival change that
reading. In a pure-birth reference population the doubling time is $\log2/b$, so
$T=20$ equals $2/\log2\approx2.885$ doubling times --- which is again not a
population doubling count in the full model with death, state changes and
selection. A physical division rate $b_{\rm phys}$ would map one pulse to
$1/b_{\rm phys}$ physical time units, but choosing a time unit cannot by itself
establish that the regimen is acute rather than protracted. Finally, the
objective extends well beyond the course: the finite-clearing construction of
Proposition~\ref{prop:clearing} adds another $152$ time units before the
alternate source begins, so $20$ is not the horizon over which eventual rescue is
observed.

Likewise, a per-division probability of reaching a coarse-grained durable state
is not comparable with a per-base DNA mutation rate. A meaningful comparison
would require a specified target size, a route to resistance, a heritability
assumption and an acquisition mechanism. The values $.01$ and $.04$ are synthetic
and would need independent justification in any application; calling them
phenotypic switching does not supply that justification, particularly as $M$ here
is irreversible and autonomous.

The practical content is about what a model must preserve, not about dosing. Our
worked example fixes $AA$ founders and compares exactly two equal-duration orders
under a declared continuation. To assess an instance of this phenomenon one would
need the joint sister-state distribution conditional on maternal state,
type-specific division and death hazards, switching rates under intervention,
division-associated acquisition, the preparation law, and the survival and
reversion behaviour of the durable state. Marginal daughter frequencies and total
population means, however well validated, cannot identify the dependence that
decides this comparison.

\paragraph{Remaining limits.} The result is confined to the literal source, the
stipulated acquisition convention, the $AA$ preparation, the two pulse orders,
the certified killing interval and the declared continuations. The effect is
small. Feedback feasibility is bounded but not optimised. Finite clearing is
expensive. Acquisition rates and calendar times are not calibrated. Competition,
spatial structure, immune effects, pharmacokinetics, toxicity, density
dependence, inherited mutant phenotypes, simultaneous continuation perturbations
and global schedule optimisation are all absent. None of these is needed for the
counterexample, and none of them repairs the failure it exhibits: exact agreement
of the complete mean dynamics does not determine the preferred order for eventual
extinction.

\appendix

\section{Mechanism details}\label{sec:mechdetail}

\subsection{Collecting the expansion}
Substituting $x=\one-\eps h+\eps^2k+O(\eps^3)$, $\mu_i=\eps\kappa_i$ and $y=\rho$
into \eqref{eq:field} and using $Q_e\one=0$:
\begin{align*}
 (Q_ex)_i&=-\eps(Q_eh)_i+\eps^2(Q_ek)_i,\\
 d_i(1-x_i)&=\eps d_ih_i-\eps^2d_ik_i,\\
 (1-\mu_i)D_i(x,x)&=1-\eps\kappa_i-2\eps(\Lambda h)_i
   +\eps^2\bigl\{2(\Lambda k)_i+D_i(h,h)+2\kappa_i(\Lambda h)_i\bigr\}+O(\eps^3),\\
 \mu_i\rho(\Lambda x)_i&=\eps\kappa_i\rho-\eps^2\kappa_i\rho(\Lambda h)_i+O(\eps^3),\\
 -x_i&=-1+\eps h_i-\eps^2k_i.
\end{align*}
The left side is $x'=-\eps h'+\eps^2k'$. Matching $O(\eps)$ gives
$-h'=-(Q_eh)_i+d_ih_i+b\{-2(\Lambda h)_i+\kappa_i\rho-\kappa_i+h_i\}$, that is
\eqref{eq:h}. Matching $O(\eps^2)$ gives
$k'=(Q_ek)_i-d_ik_i+b\{2(\Lambda k)_i+D_i(h,h)+(2-\rho)\kappa_i(\Lambda h)_i-k_i\}$,
that is \eqref{eq:k}. The $O(\eps^2)$ acquisition contribution
$b(2-\rho)\kappa_i(\Lambda h)_i$ involves the joint law only through $\Lambda$ and
therefore cancels from $\Delta=k_I-k_J$.

\subsection{The quasi-monotonicity inequalities}
The order \eqref{eq:order} has least element $RR$, greatest element $AA$, six
covering pairs and seven nonempty up-sets, namely
\begin{gather*}
\{AA\},\quad \{AU,AA\},\quad \{UU,AU,AA\},\quad \{AR,AU,AA\},\\
\{UU,AR,AU,AA\},\quad \{UU,UR,AR,AU,AA\},\quad \text{and all six types}.
\end{gather*}
For each of the six covering pairs $i\lessdot j$ and each up-set indicator $\chi$
with $\chi_i=\chi_j$, Lemma~\ref{lem:cone} requires
$\sum_k(A_{jk}-A_{ik})\chi_k\ge0$. This is $32$ rational inequalities per phase.
Exact evaluation confirms all $64$; the minimum value attained is $0$, at the
pair $UR\lessdot UU$ against the ray $\chi=\one_{\{AA\}}$, in both phases. The
script \texttt{verify\_claims.py} in the source package re-runs the enumeration
and the arithmetic.

\subsection{Why the asymptotic remainder is not usable}
Let $P,V_2,V_3$ bound the first three $\eps$-derivatives of $x$ in sup norm. With
logarithmic Jacobian norm at most $.09$ and the elementary bounds
$|f_\eps|\le.04$, $|f_{x\eps}|\le.12$, $|f_{xx}|\le.2$ and $|f_{xx\eps}|\le.08$,
the variational equations and positive scalar comparison give
\[
 P'=.09P+.04,\qquad V_2'=.09V_2+.2P^2+.24P,\qquad
 V_3'=.09V_3+.6PV_2+.36V_2+.24P^2,
\]
all starting from zero, whence
$\|x_I-x_J-\eps^2\Delta\|\le V_3(T)\eps^3/3$. At $T=20$ the coarse algebraic
bound $e^{1.8}<9$ yields $P\le7.2$, $V_2\le2177.28$ and $V_3<1.84\times10^6$, so
the permitted error at $\eps=.1$ exceeds $600$. A valid asymptotic statement can
therefore be entirely useless for a $10^{-6}$ decision. This motivated both the
exact identity of Proposition~\ref{prop:exact} and the enclosed flows.
No larger search for a sharper remainder was undertaken.

\section{Exact enclosure method and evidence boundary}\label{sec:certificate}

\subsection{Why the flow intervals are rigorous}
Write the literal field as $d+Bz+C(z,z)$ with nonnegative quadratic
coefficients. If $c_0=z$, the normalised time derivatives obey
\begin{equation}
 (k+1)c_{k+1}=\one_{k=0}d+Bc_k+\sum_{j=0}^kC(c_j,c_{k-j}),\label{eq:jets}
\end{equation}
by differentiating the polynomial ODE. Applying \eqref{eq:jets} with $c_0$ equal
to the \emph{whole} invariant cube and outward interval arithmetic bounds the
normalised twelfth derivative everywhere in each phase. If the largest absolute
endpoint is $M$, Taylor's theorem bounds the degree-$11$ local remainder by
$M(\Delta t)^{12}$. A small final Taylor term is never used as an unproved error
estimate.

The Jacobian of the full field is Metzler on the cube, with molecular row sums at
most $b-d_i\le.09$ and mutant row sum $\beta-\delta=.08$. The maximal-coordinate
differential inequality for the difference of two solutions therefore gives
sup-norm Lipschitz factor $e^{.09t}$; the implementation uses the rational upper
bound $e^{.09\Delta t}\le1/(1-.09\Delta t)$. This estimate applies to the segment
between any two probability solutions, so it does not require a componentwise
interval box to stay inside the cube.

All interval endpoints are integers divided by $2^{110}$. Multiplication and
positive integer division round outward, and rational source coefficients are
enclosed before use. Each step evaluates a degree-$11$ Horner polynomial on
intervals, projects the midpoint into the cube, and adds the distance from that
centre to both endpoints, including the projection, to the local remainder. If
the old error is $E$, the new error satisfies
\[
 E_{\rm next}\le\frac{E}{1-.09\Delta t}+E_{\rm rounding}+M(\Delta t)^{12}.
\]
Errors carry across phase boundaries without reset. With $\Delta t=1/20$ each
schedule takes $400$ steps, and each central order-gap enclosure has width below
$4\times10^{-16}$, more than nine orders of magnitude smaller than the margin.
This is exact validated arithmetic under the foregoing conventional soundness
proof; it is not a machine-verified integrator.

\subsection{Scalar and linear consequences}\label{sec:linear}
For \eqref{eq:concretefloor} the interval $[0,.8]$ is invariant for $a$ in
\eqref{eq:riccati}, whose field derivative $-.3228-.198a$ is negative there, so
the flow is $1$-Lipschitz and the same normalised-derivative construction --- one
coordinate, step $.1$, degree $11$, $200$ steps --- accumulates local and rounding
error without growth. It encloses $a(20)$ in
$(.002472578295567906,.002472578295567911)$. The closed-form Riccati solution was
evaluated as a diagnostic and matches to floating precision; it is not the
certificate.

For the linear rewards of Section~\ref{sec:continuation}, the exact source
matrices satisfy $\|A\|_\infty\le3$ with row sums at most $.09$. The global bounds
$9\|v\|_\infty<1000$ and $9\|w\|_\infty<1000$ control the reference rewards, and
at each step the centre is checked to have norm at most $950$, so the local
reference trajectory over $\Delta t=.05$ stays below $1000$. A degree-$15$ Taylor
evaluation then has local remainder at most $1000\cdot3^{16}(.05)^{16}/16!$.
Outward arithmetic, explicit centre rounding and the same rational logarithmic
norm growth enclose the course means, so the continuation constants $U_s$ and
$C_s$ do not rely on an unenclosed floating-point matrix exponential.

Regret, founder and perturbation comparisons are rational operations on the
preserved central certificate. The interval theorem of
Lemma~\ref{lem:curvature} adds only four endpoint flows; it is not a validated
high-dimensional sensitivity search. Floating-point roots and curves are
illustrations and are inputs to no certified sign conclusion.

\section{Evidence map and reproduction}\label{sec:evidencemap}

\begin{center}\small
\begin{tabular}{lp{.72\linewidth}}\toprule
Evidence&Claims supported\\\midrule
C: conventional&Branching construction and probability bridge; extinction
endpoint; expansion and exact response identities; covariance sign via Nagumo and
Harris; curvature bound and the secant argument; feedback envelope; occupation,
coupling and clearing bounds.\\
E: exact arithmetic&Mean-operator identity; clearing weight; the $64$
quasi-monotonicity inequalities; central and endpoint flow intervals; scalar
floor; linear reward bounds $U_s,C_s$; regret enclosure, founder increment signs
and the two continuation radii.\\
K: Lean kernel&Finite source normalisation and augmented linearisation; the
offspring expansion and covariance-forcing algebra; continuation weight;
scalar-envelope and barrier algebra; regret inequalities; decision-preservation
and presentation margins. Not the complete stochastic theorem.\\
N: diagnostics&Table~\ref{tab:events} probabilities; the smooth gap curves and
their zero crossings; the leading-order transport curves of
Figure~\ref{fig:decision}(c). No certified conclusion rests on these.\\
\bottomrule
\end{tabular}
\end{center}

The probability construction, ODE correspondence, continuation identification,
mechanism, curvature argument and feedback theorem are conventional proofs. The
numerical claims use exact outward arithmetic under Appendix~\ref{sec:certificate}.
Lean modules under \texttt{proofs/EvolutionaryRescue/} check the literal source
linearisation, the continuation weight, the scalar-envelope algebra, the regret
inequalities and the presentation margins; strict verification treats compiler
warnings as errors. The thin \texttt{Main.lean} imports those modules. It does
\emph{not} formalise the stochastic process, the interval integrator or the
complete probability theorem, and no end-to-end Lean formalisation is claimed.
Exact rational arithmetic does not by itself make the integrator kernel-verified;
that boundary is stated here rather than buried in a repository note.

The source package accompanying this manuscript contains \texttt{main.tex}, the
figure sources and two scripts: \texttt{figures.py}, which rebuilds both figures
from the preserved certificates, and \texttt{verify\_claims.py}, which replays
every printed numerical claim, asserting the exact ones and tagging the
diagnostics. The preserved original certificate is
\texttt{results/decision\_certificate.json}; the additional endpoint flows,
continuation constants, scalar floor and decision consequences are under
\texttt{results/postproof/}. The campaign workspace
\texttt{REPRODUCE.md} gives a bounded command that audits the original
certificate, rebuilds the selected consequences and figures, strictly verifies
the Lean root and compiles this manuscript; the original four-flow producer
remains available for independent replay. Reported reproduction receipts record
provenance, not mathematical proof.

\begin{thebibliography}{9}

\bibitem{gunnarsson2020}
E.~B. Gunnarsson, S.~De, K.~Leder and J.~Foo.
Understanding the role of phenotypic switching in cancer drug resistance.
\emph{Journal of Theoretical Biology} \textbf{490} (2020), 110162.
\url{https://arxiv.org/abs/2001.07690}.

\bibitem{gunnarsson2025}
E.~B. Gunnarsson, B.~V. Magn\'usson and J.~Foo.
Optimal dosing of anti-cancer treatment under drug-induced plasticity.
arXiv:2412.16391v2 (2025). \url{https://arxiv.org/abs/2412.16391}.

\bibitem{harris1960}
T.~E. Harris.
A lower bound for the critical probability in a certain percolation process.
\emph{Mathematical Proceedings of the Cambridge Philosophical Society}
\textbf{56} (1960), 13--20.

\bibitem{iwasa2006}
Y.~Iwasa, M.~A. Nowak and F.~Michor.
Evolution of resistance during clonal expansion.
\emph{Genetics} \textbf{172} (2006), 2557--2566.
\url{https://doi.org/10.1534/genetics.105.049791}.

\bibitem{orr2014}
H.~A. Orr and R.~L. Unckless.
The population genetics of evolutionary rescue.
\emph{PLOS Genetics} \textbf{10}(8) (2014), e1004551.
\url{https://doi.org/10.1371/journal.pgen.1004551}.

\bibitem{patil2026}
S.~Patil, A.~Ahmed, Y.~Viossat and R.~Noble.
Preventing evolutionary rescue in cancer using two-strike therapy.
\emph{Genetics} \textbf{232}(2) (2026), iyaf255.
\url{https://doi.org/10.1093/genetics/iyaf255}.

\end{thebibliography}
\end{document}
